% 第1章第6节：本章小结。

\section{本章小结}
\label{sec:dds-summary}

章首从一分钟一步的温度递推出发，问了几个问题：轨迹停在哪里，偏差会消失还是放大，非线性更新的长期结果是否同样简单，计算机算出的数字有多可信。回答这些问题时，最容易出错的地方是把不同的结论混为一谈。表~\ref{tab:dds-summary}把本章的主要结论与各自所需的条件并列：递推在每一步有定义，只保证有限轨迹可以算出；平衡点存在，只说明停在那里不会动；稳定、渐近稳定与指数稳定各有不同的量词；长期收敛还可能被瞬态放大、周期轨道或敏感依赖所取代。这些结论都不能互相替代。

\begin{table*}[htbp]
  \centering
  \begin{tabularx}{\linewidth}{>{\raggedright\arraybackslash}p{30mm}>{\raggedright\arraybackslash}X>{\raggedright\arraybackslash}p{42mm}}
    \toprule
    结论 & 所需条件 & 本章依据\\
    \midrule
    轨迹每步有定义 & 更新把状态空间映回自身，或初值位于正向不变集 & 定义~\ref{def:dds-state-model}、第~\ref{sec:dds-orbits-invariant}小节\\
    线性轨迹的显式解 & 线性或仿射更新；时变时保留乘积次序 & 定理~\ref{thm:dds-variation-constants}、式~\eqref{eq:dds-time-varying-solution}\\
    平衡点或周期轨道存在 & 解$F(\bm x)=\bm x$或$F^p(x)=x$并核查最小周期 & 第~\ref{sec:dds-equilibria-periodic}小节\\
    线性系统渐近（指数）稳定 & $\rho(\bm A)<1$；谱半径为$1$时另查单位圆上的Jordan结构 & 定理~\ref{thm:dds-spectral-stability}\\
    非线性平衡点局部稳定或不稳定 & 在平衡点可微，Jacobian矩阵的谱不在单位圆上 & 定理~\ref{thm:dds-linearization}\\
    有限步内不放大 & 谱判据不保证；需估计$\lVert\bm A^n\rVert$或改用Lyapunov范数 & 图~\ref{fig:dds-transient-growth}、定理~\ref{thm:dds-discrete-lyapunov-equation}\\
    区域内的吸引 & 紧子水平集上$\Delta V<0$；非严格下降时检查零下降集合中的不变部分 & 定理~\ref{thm:dds-sublevel-basin}、定理~\ref{thm:dds-lasalle}\\
    极限集合非空且被趋近 & 映射连续、轨道包含在紧集中 & 定理~\ref{thm:dds-omega-limit-properties}\\
    计算误差有界 & 收缩且计算序列留在适用区域；否则误差可按$L^n$增长 & 定理~\ref{thm:dds-error-propagation}\\
    持续输入偏差的影响 & 稳定时有界但一般不消失；衰减扰动的影响趋零 & 定理~\ref{thm:dds-input-response}\\
    \bottomrule
  \end{tabularx}
  \caption{本章主要结论与所需条件。左列为要得到的结论，中列为本章使用的充分条件或需要额外检查的情形，右列为相应的定义、定理或图。}
  \label{tab:dds-summary}
\end{table*}

线性更新与非线性更新提供了两种不同深度的认识。线性递推可以显式写出：初态经矩阵幂传播，每一步的输入按“注入时刻、剩余步数”累计，特征值决定各方向是增长、衰减还是旋转，状态转移矩阵记录时变情形的先后次序，结合律又把串行递推变成可并行的扫描。非线性递推一般没有显式解，我们得到的是局部或区域性的认识：在平衡点附近借线性化判断，在更大的区域里借Lyapunov函数的子水平集证明吸引，在更长的时间尺度上用极限集合、分岔与Lyapunov指数描述轨道的去向及其对参数和初值的依赖。$r=4$的Logistic映射说明，确定的规则同样可以产生无法逐点长期预测的轨道，而第~\ref{sec:dds-error-propagation}节的误差传播界则告诉我们，在什么条件下数值计算可以信任、在什么条件下误差会被系统本身放大。

离散模型逐步规定状态：只要知道这一步，下一步就确定了。很多物理和学习过程却是用瞬时变化率描述的，例如温度的变化速度与当前温差成正比，参数沿负梯度方向连续流动。这时需要重新回答本章的每个问题：给定初值的解是否存在、是否唯一，怎样求解，平衡点何时稳定，以及用离散步长计算时误差如何累计。第~\ref{chap:ordinary-differential-equations}章在连续时间中重建这些概念，并说明本章的递推怎样作为数值方法回到连续问题中。再往后，第~\ref{chap:discrete-time-stochastic-processes}章在递推中加入随机噪声和历史信息，研究带噪的迭代如何转移、累计，以及长期分布与时间平均怎样变化。
